%% LyX 2.4.4 created this file.  For more info, see https://www.lyx.org/.
%% Do not edit unless you really know what you are doing.
\documentclass[english]{article}
\usepackage[T1]{fontenc}
\usepackage[latin9]{inputenc}
\usepackage{enumitem}
\usepackage{amsmath}
\usepackage{amsthm}
\usepackage{cancel}
\usepackage{geometry}
\geometry{verbose,tmargin=1in,bmargin=1in,lmargin=1in,rmargin=1in}

\makeatletter
%%%%%%%%%%%%%%%%%%%%%%%%%%%%%% Textclass specific LaTeX commands.
\numberwithin{equation}{section}
\numberwithin{figure}{section}
\newlength{\lyxlabelwidth}      % auxiliary length 

%%%%%%%%%%%%%%%%%%%%%%%%%%%%%% User specified LaTeX commands.
\usepackage{fancyhdr}
\usepackage{lastpage}
\pagestyle{fancy}
\usepackage{enumitem}
\usepackage{ifthen}
%sets math fonts to be more readable
\everymath=\expandafter{\the\everymath\displaystyle}
%\DeclareMathSizes{10}{10}{10}{10}
\usepackage{multicol}

\usepackage{tikz}
\usetikzlibrary{plotmarks}
\usepackage{pgfplots}
\pgfplotsset{compat=newest}

\usepackage{calc}
\usepackage[nomessages]{fp}

\usepackage{advdate}    % Advancing/saving dates
\usepackage{datetime}   % Dates formatting
\usepackage{datenumber} % Counters for dates
% Added by lyx2lyx
\setlength{\parskip}{\smallskipamount}
\setlength{\parindent}{0pt}

\makeatother

\usepackage{babel}
\begin{document}
\renewcommand{\author}{Dr.\,James Ashe}

\newcommand{\classNum}{137}

\newcommand{\classSec}{B}

\newcommand{\nameType}{Name:}

\newcommand{\examType}{Exam 4}

\renewcommand{\date}{\formatdate{6}{5}{2013}}

%%First page's header%%%%%%%%%%%%%%%%%%%%%%
\fancypagestyle{firststyle} %%header settings for first pagr
{
	\fancyhead[L]{MTH \classNum{}(\classSec{}) \author{}\\
	\textbf{\examType{}} \date{}}
	\fancyhead[C]{\textbf{\nameType}}
	\setlength{\headsep}{14pt}
}
%%%%%%%%%%%%%%%%%%%%%%%%%%%%%%%%%%%%%%%%%%

%%Next page's header%%%%%%%%%%%%%%%%%%%%%%%
\setlist{nolistsep}
\lhead{\classNum{}(\classSec{})} 
\chead{\textbf{\nameType}} 
\cfoot{\thepage{}~of~\pageref{LastPage}} 
\setlength{\headsep}{5pt} 
%%%%%%%%%%%%%%%%%%%%%%%%%%%%%%%%%%%%%%%%%%%

%%Various settings; see the preamble for other settings%%%%%
%%\setenumerate[0]{leftmargin=12pt,itemindent=0pt,labelwidth=0pt}
\renewcommand{\labelenumi}{\textbf{\arabic{enumi}.}}
\renewcommand{\labelenumii}{\textbf{\alph{enumii}.}}
\thispagestyle{firststyle}
%%%%%%%%%%%%%%%%%%%%%%%%%%%%%%%%%%%%%%%%%%
\newcommand{\xmax}{10}
\newcommand{\xmin}{-10}
\newcommand{\xstep}{0} %must be xmin+n
\newcommand{\ymax}{10}
\newcommand{\ymin}{-10}
\newcommand{\ystep}{0} %must be ymin+n

\newcommand{\Dspace}{\vspace{.75cm}}

\textbf{You must show your work to receive credit. }Unless stated
otherwise, all problems are worth 5 points. 
\begin{enumerate}
\item Evaluate the following {[}4 pts each{]}:

\begin{enumerate}
\item $-2^{-3}$\Dspace
\item $\log_{\pi}1$\Dspace
\item $\ln\left(e^{3}\right)$\Dspace
\item ${\displaystyle \lim_{x\rightarrow-\infty}7^{x}}$\Dspace
\item ${\displaystyle \lim_{x\rightarrow\infty}\log_{7}x}$\Dspace
\end{enumerate}
\end{enumerate}
\begin{enumerate}[resume]
\item Let $f(x)=2^{x}-3$.\begin{multicols}{2}

\begin{enumerate}
\item Graph $f(x)$.\Dspace
\item Find the domain and range of $f(x)$.\vfill\columnbreak
\item []\renewcommand{\xmax}{10}
\renewcommand{\xmin}{-10}
\renewcommand{\xstep}{0} %must be xmin+n
\renewcommand{\ymax}{10}
\renewcommand{\ymin}{-10}
\renewcommand{\ystep}{0} %must be ymin+n
%%%%%%%
\begin{tikzpicture} 
\begin{axis}[height=5cm, 
	width=5cm, 
	axis lines=middle,
	minor x tick num={4},
	minor y tick num={4},
	grid=both,
	%xtick={0,50,...,250},
	%ytick={0,10,20,...,40},
	%axis line style={-latex}, 
	xmin=\xmin,xmax=\xmax, 
	ymin=\ymin,ymax=\ymax,
	%restrict y to domain=-1:10,
	no markers, 
	xlabel={$x$},ylabel={$y$}, 
	%every axis x label/.append style={anchor=north}, 
	%every axis y label/.append style={anchor=east}
	]
	
	%%\addplot[color=red,solid,mark=*] coordinates {(-3,-4)};
	%%\addplot[domain=\xmin:\xmax,thick,samples=100,draw=red,<->]{-(x+2)*(x-2)*(x-5)} node[pos=.65,right] {$f(x)$};
	
\end{axis}
\end{tikzpicture}
\end{enumerate}
\end{multicols}
\item Let $g(x)=\log_{4}\left(x-1\right)$.\begin{multicols}{2}

\begin{enumerate}
\item Graph $g(x)$.\Dspace
\item Find the domain and range of $g(x)$.\vfill\columnbreak
\item []\renewcommand{\xmax}{10}
\renewcommand{\xmin}{-10}
\renewcommand{\xstep}{0} %must be xmin+n
\renewcommand{\ymax}{10}
\renewcommand{\ymin}{-10}
\renewcommand{\ystep}{0} %must be ymin+n
%%%%%%%
\begin{tikzpicture} 
\begin{axis}[height=5cm, 
	width=5cm, 
	axis lines=middle,
	minor x tick num={4},
	minor y tick num={4},
	grid=both,
	%xtick={0,50,...,250},
	%ytick={0,10,20,...,40},
	%axis line style={-latex}, 
	xmin=\xmin,xmax=\xmax, 
	ymin=\ymin,ymax=\ymax,
	%restrict y to domain=-1:10,
	no markers, 
	xlabel={$x$},ylabel={$y$}, 
	%every axis x label/.append style={anchor=north}, 
	%every axis y label/.append style={anchor=east}
	]
	
	%%\addplot[color=red,solid,mark=*] coordinates {(-3,-4)};
	%%\addplot[domain=\xmin:\xmax,thick,samples=100,draw=red,<->]{-(x+2)*(x-2)*(x-5)} node[pos=.65,right] {$f(x)$};
	
\end{axis}
\end{tikzpicture}
\end{enumerate}
\end{multicols}
\item Use the properties of logarithms to expand the expression as much
as possible.

\begin{enumerate}
\item $\log_{7}\left(3y\right)$\vfill
\item $\log\left(\frac{10^{6}}{x}\right)$\vfill\newpage
\end{enumerate}
\item Use the properties of logarithms to rewrite the expression as a single
logarithm.

\begin{enumerate}
\item $\log_{c}\left(16\right)-\log_{c}\left(2\right)$\vfill
\item $\ln(2)+6\ln\left(x\right)$\vfill
\end{enumerate}
\item Solve for $x$. Write your answer as an exact and approximate value
rounded to three decimal places {[}4 pts each{]}.

\begin{enumerate}
\item $4^{x}=3$\vfill
\item $e^{2x+1}=10$\vfill
\item $\log_{6}x=2$\vfill
\item $\ln(40)=3-\ln(x)$\vfill
\item $\log_{2}\left(x+3\right)+\log_{2}\left(x-3\right)=3$\vfill\Dspace\newpage
\end{enumerate}
\item If \$500 is invested in an account at $5\%$ annual interest compounded
monthly, what is the amount in the account after 9 years?\vfill\Dspace
\item How long will it take \$500 to grow to \$1500 if it is invested at
10\% compounded continuously?\vfill\Dspace
\item Find and graph the focus, directrix, and vertex of the parabola $y=\frac{1}{8}(x-4)^{2}+1$.\\
\renewcommand{\ymax}{10}
\renewcommand{\ymin}{-5}
\renewcommand{\xmax}{14}
\renewcommand{\xmin}{-6}
%%%%%%%
\begin{tikzpicture} 
\begin{axis}[height=6cm, 
	width=6cm, 
	axis lines=middle,
	minor x tick num={4},
	minor y tick num={4},
	%xtick={0,50,...,250},
	%ytick={0,10,20,...,40},
	%grid=both,
	%axis line style={-latex}, 
	xmin=\xmin,xmax=\xmax, 
	ymin=\ymin,ymax=\ymax,
	%restrict y to domain=-4:6,
	%no markers, 
	xlabel={$x$},ylabel={$y$}, 
	%every axis x label/.append style={anchor=north}, 
	%every axis y label/.append style={anchor=east}
	]
	
	%\addplot[color=black,solid,mark=*] coordinates {(4,1)};
	%\addplot[color=black,solid,mark=*] coordinates {(4,3)} node[above,xshift=0pt] {focus};
	%\addplot[domain=\xmin:\xmax,thick,samples=100,draw=blue,dashed]{-1} node[below,xshift=-1cm] {directrix:} node[below,xshift=-1cm,yshift=-.5cm] {$y=-1$};
	\addplot[domain=\xmin:\xmax,thick,samples=100,draw=red,<->]{1/8*(x-4)^2+1} node[left] {$f(x)$};
	
\end{axis}
\end{tikzpicture}\Dspace\Dspace
\item Find and graph the equation of the parabola with focus $(-4,-5)$
and directrix $y=-3$. ($p=-1$)\\
\renewcommand{\ymax}{6}
\renewcommand{\ymin}{-14}
\renewcommand{\xmax}{6}
\renewcommand{\xmin}{-14}
%%%%%%%
\begin{tikzpicture} 
\begin{axis}[height=6cm, 
	width=6cm, 
	axis lines=middle,
	minor x tick num={4},
	minor y tick num={4},
	%xtick={0,50,...,250},
	%ytick={0,10,20,...,40},
	%grid=both,
	%axis line style={-latex}, 
	xmin=\xmin,xmax=\xmax, 
	ymin=\ymin,ymax=\ymax,
	%restrict y to domain=-1:10,
	%no markers, 
	xlabel={$x$},ylabel={$y$}, 
	%every axis x label/.append style={anchor=north}, 
	%every axis y label/.append style={anchor=east}
	]
	
	%\addplot[color=black,solid,mark=*] coordinates {(-4,-4)};
	%\addplot[color=black,solid,mark=*] coordinates {(-4,-5)} node[below] {focus};
	%\addplot[domain=\xmin:\xmax,thick,samples=100,draw=blue,dashed]{-3} node[above,xshift=-1cm] {directrix:} node[below,xshift=-1cm,yshift=-.5cm] {$y=2$};
	%\addplot[domain=\xmin:\xmax,thick,samples=100,draw=red,<->]{-1/4*(x+4)^2-4} node[left] {$f(x)$};
	
\end{axis}
\end{tikzpicture}\Dspace\Dspace\newpage
\end{enumerate}

\end{document}
